\documentclass[dvisvgm,tikz,border=3pt]{standalone}
\usepackage{amsmath,amssymb}
\usepackage{pgfplots}
\usepackage{CJKutf8}
\pgfplotsset{compat=1.18}
\usetikzlibrary{arrows.meta,calc,positioning}

\definecolor{themebg}{HTML}{2e362c}
\definecolor{themefg}{HTML}{edf4e8}
\definecolor{themegray}{HTML}{9ea897}
\definecolor{themecurve}{HTML}{7eb6ff}
\definecolor{themealert}{HTML}{ff7b7b}
\definecolor{themeamber}{HTML}{f5b942}

\pgfdeclarelayer{background}
\pgfdeclarelayer{foreground}
\pgfsetlayers{background,main,foreground}

\begin{document}
\begin{CJK*}{UTF8}{gbsn}
\pagecolor{themebg}
\begin{tikzpicture}[color=themefg, text=themefg]

\begin{axis}[
    width=13.5cm,
    height=8.2cm,
    axis lines=middle,
    axis line style={color=themefg, -{Stealth}},
    tick style={color=themefg},
    tick label style={color=themefg, font=\scriptsize},
    every axis label/.style={color=themefg},
    xmin=-2.2, xmax=2.2,
    ymin=-0.5, ymax=4.2,
    xtick={-2,-1,0,1,2},
    ytick={1,2,3,4},
    clip=false,
    xlabel={$x$},
    ylabel={$y$},
    title style={font=\footnotesize\bfseries, at={(0.5,1.06)}, anchor=south},
    title={一元局部模型母图：基点 $x=0$ 处的逐级 Taylor 多项式逼近与余项收敛},
]

% 1. 原真实函数 f(x) = e^x
\addplot[themecurve, line width=2.4pt, domain=-2.1:1.4, samples=120] {exp(x)};
\node[text=themecurve, font=\small\bfseries, fill=themebg, inner sep=1pt, anchor=west] at (axis cs:1.42,3.9) {$y=e^x$};

% 2. 零阶模型 T0(x) = 1
\addplot[themegray, dashed, line width=1.2pt, domain=-2.0:2.0] {1};
\node[text=themegray, font=\scriptsize, fill=themebg, inner sep=1pt, anchor=south east] at (axis cs:-1.2,1.05) {$T_0(x)=1$};

% 3. 一阶模型 T1(x) = 1 + x
\addplot[themeamber, line width=1.6pt, domain=-1.5:2.0] {1 + x};
\node[text=themeamber, font=\scriptsize, fill=themebg, inner sep=1pt, anchor=north west] at (axis cs:1.8,2.7) {$T_1(x)=1+x$ (切线)};

% 4. 二阶模型 T2(x) = 1 + x + x^2/2
\addplot[themealert, dashdotted, line width=1.6pt, domain=-2.0:1.6, samples=100] {1 + x + 0.5*x^2};
\node[text=themealert, font=\scriptsize, fill=themebg, inner sep=1pt, anchor=south west] at (axis cs:1.4,3.1) {$T_2(x)=1+x+\frac{x^2}{2}$ (密切抛物线)};

% 5. 三阶模型 T3(x) = 1 + x + x^2/2 + x^3/6
\addplot[themefg!90, densely dashed, line width=1.4pt, domain=-2.1:1.5, samples=100] {1 + x + 0.5*x^2 + (1/6)*x^3};
\node[text=themefg!90, font=\scriptsize, fill=themebg, inner sep=1pt, anchor=north west] at (axis cs:-1.6,0.35) {$T_3(x)=1+x+\frac{x^2}{2}+\frac{x^3}{6}$};

\begin{pgfonlayer}{foreground}
    % 基点 (0,1)
    \addplot[only marks, mark=*, mark size=3pt, fill=themecurve, draw=themefg] coordinates {(0,1)};
    \node[text=themecurve, font=\scriptsize, fill=themebg, inner sep=1pt, anchor=south east] at (axis cs:-0.08,1.15) {基点 $(0,1)$};

    % 局部模型升级标注
    \node[text=themefg, font=\scriptsize, fill=themebg, inner sep=1.2pt, anchor=north west] at (axis cs:-2.1,3.8) {
        \begin{tabular}{l}
        \textbf{局部误差阶逐级递减}：\\
        $f(x) - T_0(x) = O(x)$\\
        $f(x) - T_1(x) = o(x) = O(x^2)$\\
        $f(x) - T_2(x) = o(x^2) = O(x^3)$\\
        $f(x) - T_3(x) = o(x^3) = O(x^4)$
        \end{tabular}
    };
\end{pgfonlayer}

\end{axis}
\end{tikzpicture}
\end{CJK*}
\end{document}
